\begin{tikzpicture}[
  font=\sffamily,
  pillar/.style={
    rectangle, rounded corners=4pt, line width=1pt,
    text width=4.0cm, align=left, inner sep=7pt, minimum height=4.4cm, font=\sffamily\scriptsize
  }
]
  \node[font=\sffamily\bfseries\small, ceindigo] at (4.5, 2.85)
    {PATRÓN DE DESAGREGACIÓN: TRES NEGOCIOS BAJO UN MISMO TECHO (HAGEL Y SINGER)};

  \node[pillar, draw=ceindigo, fill=cepanel] (p1) at (0, 0) {
    \textcolor{ceindigo}{\bfseries\footnotesize 1. GESTIÓN DE INFRAESTRUCTURA}\\[4pt]
    \textbf{Lógica económica:} altas economías de escala; el coste unitario cae con el volumen.\\[4pt]
    \textbf{Cultura:} estandarización, eficiencia y minimización de costes.\\[4pt]
    \textbf{Competición:} por tamaño operativo.\\[4pt]
    \textit{Ej.: torres de telecomunicación, centros de datos AWS, red logística.}
  };

  \node[pillar, draw=cecrimson, fill=cepanelcrimson] (p2) at (4.5, 0) {
    \textcolor{cecrimson}{\bfseries\footnotesize 2. INNOVACIÓN DE PRODUCTO}\\[4pt]
    \textbf{Lógica económica:} velocidad de llegada al mercado; prima por novedad temprana.\\[4pt]
    \textbf{Cultura:} foco en el talento, experimentación y tolerancia al fallo.\\[4pt]
    \textbf{Competición:} por talento y propiedad intelectual.\\[4pt]
    \textit{Ej.: I+D farmacéutica, estudios de diseño, desarrolladores de apps.}
  };

  \node[pillar, draw=ceteal, fill=cepanelteal] (p3) at (9.0, 0) {
    \textcolor{ceteal!90!black}{\bfseries\footnotesize 3. RELACIÓN CON CLIENTES}\\[4pt]
    \textbf{Lógica económica:} economías de alcance; amortizar el coste de captación con más productos.\\[4pt]
    \textbf{Cultura:} orientación al servicio, confianza y personalización.\\[4pt]
    \textbf{Competición:} por cuota de bolsillo.\\[4pt]
    \textit{Ej.: banca privada, operadoras comerciales, plataformas de asesoría.}
  };

  \draw[<->, >=Stealth, dashed, thick, cecrimson] (p1.east) -- (p2.west);
  \draw[<->, >=Stealth, dashed, thick, cecrimson] (p2.east) -- (p3.west);

  \node[font=\sffamily\itshape\scriptsize, ceslate] at (4.5, -2.7)
    {Regla de Osterwalder: forzar estas tres culturas opuestas en una única estructura crea fricción crónica; conviene separarlas o asociarse.};
\end{tikzpicture}